\begin{derivation}[从\eqnrefs{eq:sm-left-weyl-anomaly-set-dp11}到\eqnrefs{eq:sm-gauge-anomaly-cancellation-dp32}]
先说明为什么规范反常只需要检查正文列出的那些群论迹。对一个左手 Weyl 费米子，三个外部规范玻色子与费米子环相连时，两种相反顶角次序的群因子分别为
\begin{equation*}
 \Tr_R(T^aT^bT^c),
 \qquad
 \Tr_R(T^aT^cT^b).
\end{equation*}
三角图的宇称奇部分在交换后两个外腿时是对称的，因此两种次序相加后只保留
\begin{equation*}
 \Tr_R\!\left[T^a\{T^b,T^c\}\right]
 =\frac12d_R^{abc}.
\end{equation*}
动量积分与正则化决定共同的数值前因子，而一个给定费米子谱是否无规范反常只取决于全部左手 Weyl 场的这些群因子之和是否为零。若其中一个或多个外腿属于 $U(1)_Y$，相应生成元就是超荷乘单位矩阵，于是四类局域反常系数分别化为
\begin{align*}
 [SU(N)]^2U(1)_Y &: \sum_f Y_fT(R_f),\\
 [U(1)_Y]^3 &: \sum_f d(R_f)Y_f^3,\\
 {\rm grav}^2U(1)_Y &: \sum_f d(R_f)Y_f,\\
 [SU(N)]^3 &: \sum_f d_{R_f}^{abc}.
\end{align*}
这里 $d(R_f)$ 表示其余内部指标的多重度。规范理论要求这些群论系数在全部左手 Weyl 场上求和后同时为零。

现在把一代标准模型全部改写成左手 Weyl 场
\begin{equation*}
 Q_L:(\mathbf3,\mathbf2)_{1/6},\quad
 u_R^c:(\bar{\mathbf3},\mathbf1)_{-2/3},\quad
 d_R^c:(\bar{\mathbf3},\mathbf1)_{1/3},\quad
 L_L:(\mathbf1,\mathbf2)_{-1/2},\quad
 e_R^c:(\mathbf1,\mathbf1)_{1}.
\end{equation*}
把共同的三角圈动量积分和规范耦合常数提出以后，定义纯群论反常系数
\begin{align*}
 \mathfrak A_{33Y}
 &\equiv\sum_{\text{左手 Weyl }f}d_2(f)Y_fT_3(R_f),\\
 \mathfrak A_{22Y}
 &\equiv\sum_f d_3(f)Y_fT_2(R_f),\\
 \mathfrak A_{YYY}
 &\equiv\sum_f d_3(f)d_2(f)Y_f^3,\\
 \mathfrak A_{{\rm grav}^2Y}
 &\equiv\sum_f d_3(f)d_2(f)Y_f.
\end{align*}
其中 $d_3,d_2$ 分别是颜色和弱表示维数。利用 $T(\mathbf3)=T(\bar{\mathbf3})=T(\mathbf2)=1/2$，一代标准模型逐项给
\begin{align*}
 \mathfrak A_{33Y}
 &=2\!\left(\frac16\right)\frac12
 +\left(-\frac23\right)\frac12
 +\left(\frac13\right)\frac12=0,\\
 \mathfrak A_{22Y}
 &=3\!\left(\frac16\right)\frac12
 +\left(-\frac12\right)\frac12=0,\\
 \mathfrak A_{YYY}
 &=6\left(\frac16\right)^3
 +3\left(-\frac23\right)^3
 +3\left(\frac13\right)^3
 +2\left(-\frac12\right)^3+1^3=0,\\
 \mathfrak A_{{\rm grav}^2Y}
 &=6\left(\frac16\right)
 +3\left(-\frac23\right)
 +3\left(\frac13\right)
 +2\left(-\frac12\right)+1=0.
\end{align*}
这些等式就是四类反常的群论系数本身为零，而不是只说明它们“与某个零表达式成正比”。
对 $[SU(3)_c]^3$，基本表示与反基本表示的 $d_R^{abc}$ 符号相反，而 $Q_L$ 有两个弱分量，因此
\begin{equation*}
 2d_{\mathbf3}^{abc}+d_{\bar{\mathbf3}}^{abc}+d_{\bar{\mathbf3}}^{abc}
 =(2-1-1)d_{\mathbf3}^{abc}=0.
\end{equation*}
$SU(2)$ 基本表示是伪实表示，局域三次群不变量本身为零。因此一代费米子谱的全部局域规范反常系数均消失，即得到正文\eqnrefs{eq:sm-gauge-anomaly-cancellation-dp32}。正文另行陈述的 Witten 模二条件属于大规范变换的全局拓扑定理，不由这个微扰三角图计算推出。
\end{derivation}

\begin{derivation}[从\eqnrefs{eq:sm-hypercharge-constraint-system-dp68}到\eqnrefs{eq:sm-hypercharge-solution-dp32}]
对一般 Yukawa 耦合与超荷，三种单态的规范不变性要求

\begin{align*}
 -q-h+u&=0,\\
 -q+h+d&=0,\\
 -\ell+h+e&=0,
\end{align*}
因此
\begin{equation*}
 u=q+h,\qquad d=q-h,\qquad e=\ell-h.
\end{equation*}

$[SU(2)_L]^2U(1)_Y$ 反常只由弱双重态贡献，所以
\begin{equation*}
 3q+\ell=0,
\end{equation*}
即 $\ell=-3q$。混合引力--超荷反常给
\begin{equation*}
 6q-3u-3d+2\ell-e=0.
\end{equation*}
代入 Yukawa 关系后，左边逐项化为
\begin{align*}
 6q-3(q+h)-3(q-h)+2\ell-(\ell-h)
 &=\ell+h.
\end{align*}
所以 $\ell=-h$。与 $\ell=-3q$ 联立得到 $h=3q$。

整体超荷归一化可以吸收到 $U(1)_Y$ 耦合 $g'$ 中。采用 Higgs 超荷 $h=1/2$ 的标准约定，得到
\begin{equation*}
 q=\frac16,\qquad \ell=-\frac12.
\end{equation*}
再代回 Yukawa 关系：
\begin{equation*}
 Y_u=\frac23,\qquad
 Y_d=-\frac13,\qquad
 Y_e=-1.
\end{equation*}
这就是\eqnrefs{eq:sm-hypercharge-solution-dp32}。

三次 $U(1)_Y^3$ 条件还必须独立核验。定义
\begin{equation*}
 \mathfrak A_{YYY}(q,h)
 \equiv
 6q^3-3(q+h)^3-3(q-h)^3+2(-h)^3-(-2h)^3.
\end{equation*}
使用 $\ell=-h$、$u=q+h$、$d=q-h$、$e=-2h$ 后直接展开：
\begin{align*}
 \mathfrak A_{YYY}(q,h)
 &=6h^2(h-3q)\\
 &=0,
\end{align*}
最后一步正由已经得到的 $h=3q$ 保证，因此三次阿贝尔反常没有引入新的矛盾。
\end{derivation}
